
\documentclass{article}
\usepackage{colortbl}
\usepackage{makecell}
\usepackage{multirow}
\usepackage{supertabular}

\begin{document}

\newcounter{utterance}

\twocolumn

{ \footnotesize  \setcounter{utterance}{1}
\setlength{\tabcolsep}{0pt}
\begin{supertabular}{c@{$\;$}|p{.15\linewidth}@{}p{.15\linewidth}p{.15\linewidth}p{.15\linewidth}p{.15\linewidth}p{.15\linewidth}}

    \# & $\;$A & \multicolumn{4}{c}{Game Master} & $\;\:$B\\
    \hline 

    \theutterance \stepcounter{utterance}  

    & & \multicolumn{4}{p{0.6\linewidth}}{\cellcolor[rgb]{0.9,0.9,0.9}{%
	\makecell[{{p{\linewidth}}}]{% 
	  \tt {\tiny [A$\langle$GM]}  
	 State giocando a un gioco di negoziazione cooperativa in cui dovete accordarvi su come dividere una serie di oggetti tra voi e un altro giocatore.\\ \tt \\ \tt Regole:\\ \tt (a) A Lei e all'altro giocatore viene dato un insieme di oggetti. Ognuno di voi riceve anche una funzione di valore segreta, che rappresenta il valore di ciascun tipo di oggetto.\\ \tt (b) Si scambiano messaggi con l'altro giocatore per concordare chi si aggiudica gli oggetti. Potete inviare un massimo di 5 messaggi ciascuno, oppure terminare in anticipo facendo una proposta segreta in qualsiasi momento.\\ \tt (c) A ciascuno di voi viene chiesto di inviare una proposta segreta indicando gli oggetti che desiderate, formattata tra parentesi quadre come segue: "[Proposta: <numero> <nome oggetto>, <numero> <nome oggetto>, <...>]".\\ \tt (d) Se le vostre proposte sono complementari, cioè ci sono abbastanza oggetti per soddisfare entrambe le proposte, a ciascun giocatore viene assegnato un punteggio basato sulla somma dei valori degli oggetti ricevuti. In caso contrario, entrambi ricevono zero punti.\\ \tt (e) L'obiettivo è massimizzare la somma del proprio punteggio e di quello dell'altro giocatore.\\ \tt \\ \tt Cominciamo.\\ \tt \\ \tt L'insieme degli oggetti disponibili è:\\ \tt \\ \tt 1 patata, 1 guanto, 2 righelli, 1 fischietto, 2 schede\\ \tt \\ \tt La funzione valore segreta è:\\ \tt \\ \tt patata: 0, guanto: 1, righello: 3, fischietto: 1, scheda: 1\\ \tt \\ \tt IMPORTANTE: i vostri messaggi, a meno che non si tratti di una proposta segreta, vengono trasmessi direttamente all'altro giocatore, quindi non includete alcuna risposta alle regole o testo di annuncio del vostro messaggio. Per fare una proposta segreta, utilizzate il formato indicato. Non utilizzare delle parentesi quadre quando si comunica all'altro giocatore, altrimenti verrà interpretata come una proposta segreta.\\ \tt  
	  } 
	   } 
	   } 
	 & \\ 
 

    \theutterance \stepcounter{utterance}  

    & \multicolumn{4}{p{0.6\linewidth}}{\cellcolor[rgb]{0.8,1,0.9}{%
	\makecell[{{p{\linewidth}}}]{% 
	  \tt {\tiny [A$\rangle$GM]}  
	 Ciao! Vediamo come possiamo dividere questi oggetti nel modo più vantaggioso per entrambi. Abbiamo: 1 patata, 1 guanto, 2 righelli, 1 fischietto, 2 schede.\\ \tt \\ \tt Per trovare la divisione ottimale, dovremmo capire quali oggetti valutiamo di più. Potresti dirmi quali sono i tuoi 2-3 oggetti preferiti tra quelli disponibili? Io posso condividere le mie preferenze dopo aver sentito le tue. 
	  } 
	   } 
	   } 
	 & & \\ 
 

    \theutterance \stepcounter{utterance}  

    & & & \multicolumn{2}{p{0.3\linewidth}}{\cellcolor[rgb]{0.95,0.95,0.95}{%
	\makecell[{{p{\linewidth}}}]{% 
	  \tt {\tiny [GM$|$GM]}  
	 continue 
	  } 
	   } 
	   } 
	 & & \\ 
 

    \theutterance \stepcounter{utterance}  

    & & & \multicolumn{2}{p{0.3\linewidth}}{\cellcolor[rgb]{0.95,0.95,0.95}{%
	\makecell[{{p{\linewidth}}}]{% 
	  \tt {\tiny [GM$|$GM]}  
	 Ciao! Vediamo come possiamo dividere questi oggetti nel modo più vantaggioso per entrambi. Abbiamo: 1 patata, 1 guanto, 2 righelli, 1 fischietto, 2 schede.\\ \tt \\ \tt Per trovare la divisione ottimale, dovremmo capire quali oggetti valutiamo di più. Potresti dirmi quali sono i tuoi 2-3 oggetti preferiti tra quelli disponibili? Io posso condividere le mie preferenze dopo aver sentito le tue. 
	  } 
	   } 
	   } 
	 & & \\ 
 

    \theutterance \stepcounter{utterance}  

    & & \multicolumn{4}{p{0.6\linewidth}}{\cellcolor[rgb]{0.9,0.9,0.9}{%
	\makecell[{{p{\linewidth}}}]{% 
	  \tt {\tiny [GM$\rangle$B]}  
	 State giocando a un gioco di negoziazione cooperativa in cui dovete accordarvi su come dividere una serie di oggetti tra voi e un altro giocatore.\\ \tt \\ \tt Regole:\\ \tt (a) A Lei e all'altro giocatore viene dato un insieme di oggetti. Ognuno di voi riceve anche una funzione di valore segreta, che rappresenta il valore di ciascun tipo di oggetto.\\ \tt (b) Si scambiano messaggi con l'altro giocatore per concordare chi si aggiudica gli oggetti. Potete inviare un massimo di 5 messaggi ciascuno, oppure terminare in anticipo facendo una proposta segreta in qualsiasi momento.\\ \tt (c) A ciascuno di voi viene chiesto di inviare una proposta segreta indicando gli oggetti che desiderate, formattata tra parentesi quadre come segue: "[Proposta: <numero> <nome oggetto>, <numero> <nome oggetto>, <...>]".\\ \tt (d) Se le vostre proposte sono complementari, cioè ci sono abbastanza oggetti per soddisfare entrambe le proposte, a ciascun giocatore viene assegnato un punteggio basato sulla somma dei valori degli oggetti ricevuti. In caso contrario, entrambi ricevono zero punti.\\ \tt (e) L'obiettivo è massimizzare la somma del proprio punteggio e di quello dell'altro giocatore.\\ \tt \\ \tt Cominciamo.\\ \tt \\ \tt L'insieme degli oggetti disponibili è:\\ \tt \\ \tt 1 patata, 1 guanto, 2 righelli, 1 fischietto, 2 schede\\ \tt \\ \tt La funzione valore segreta è:\\ \tt \\ \tt patata: 7, guanto: 0, righello: 0, fischietto: 3, scheda: 0\\ \tt \\ \tt IMPORTANTE: i vostri messaggi, a meno che non si tratti di una proposta segreta, vengono trasmessi direttamente all'altro giocatore, quindi non includete alcuna risposta alle regole o testo di annuncio del vostro messaggio. Per fare una proposta segreta, utilizzate il formato indicato. Non utilizzare delle parentesi quadre quando si comunica all'altro giocatore, altrimenti verrà interpretata come una proposta segreta.\\ \tt \\ \tt \\ \tt Ciao! Vediamo come possiamo dividere questi oggetti nel modo più vantaggioso per entrambi. Abbiamo: 1 patata, 1 guanto, 2 righelli, 1 fischietto, 2 schede.\\ \tt \\ \tt Per trovare la divisione ottimale, dovremmo capire quali oggetti valutiamo di più. Potresti dirmi quali sono i tuoi 2-3 oggetti preferiti tra quelli disponibili? Io posso condividere le mie preferenze dopo aver sentito le tue. 
	  } 
	   } 
	   } 
	 & \\ 
 

    \theutterance \stepcounter{utterance}  

    & & & \multicolumn{4}{p{0.6\linewidth}}{\cellcolor[rgb]{1,0.85,0.72}{%
	\makecell[{{p{\linewidth}}}]{% 
	  \tt {\tiny [GM$\langle$B]}  
	 Ciao! Buona idea quella di condividere le nostre preferenze per trovare una divisione ottimale.\\ \tt \\ \tt I miei oggetti preferiti in ordine di preferenza sono:\\ \tt 1. La patata (la valuto molto)\\ \tt 2. Il fischietto (ha un valore moderato per me)\\ \tt 3. Gli altri oggetti (guanto, righelli, schede) hanno tutti lo stesso valore per me, abbastanza basso\\ \tt \\ \tt E tu? Quali sono i tuoi oggetti preferiti? Questo ci aiuterà a capire dove ci sono sovrapposizioni e dove invece possiamo fare scambi vantaggiosi per entrambi. 
	  } 
	   } 
	   } 
	  \\ 
 

    \theutterance \stepcounter{utterance}  

    & & & \multicolumn{2}{p{0.3\linewidth}}{\cellcolor[rgb]{0.95,0.95,0.95}{%
	\makecell[{{p{\linewidth}}}]{% 
	  \tt {\tiny [GM$|$GM]}  
	 continue 
	  } 
	   } 
	   } 
	 & & \\ 
 

    \theutterance \stepcounter{utterance}  

    & & & \multicolumn{2}{p{0.3\linewidth}}{\cellcolor[rgb]{0.95,0.95,0.95}{%
	\makecell[{{p{\linewidth}}}]{% 
	  \tt {\tiny [GM$|$GM]}  
	 Ciao! Buona idea quella di condividere le nostre preferenze per trovare una divisione ottimale.\\ \tt \\ \tt I miei oggetti preferiti in ordine di preferenza sono:\\ \tt 1. La patata (la valuto molto)\\ \tt 2. Il fischietto (ha un valore moderato per me)\\ \tt 3. Gli altri oggetti (guanto, righelli, schede) hanno tutti lo stesso valore per me, abbastanza basso\\ \tt \\ \tt E tu? Quali sono i tuoi oggetti preferiti? Questo ci aiuterà a capire dove ci sono sovrapposizioni e dove invece possiamo fare scambi vantaggiosi per entrambi. 
	  } 
	   } 
	   } 
	 & & \\ 
 

    \theutterance \stepcounter{utterance}  

    & & \multicolumn{4}{p{0.6\linewidth}}{\cellcolor[rgb]{0.9,0.9,0.9}{%
	\makecell[{{p{\linewidth}}}]{% 
	  \tt {\tiny [A$\langle$GM]}  
	 Ciao! Buona idea quella di condividere le nostre preferenze per trovare una divisione ottimale.\\ \tt \\ \tt I miei oggetti preferiti in ordine di preferenza sono:\\ \tt 1. La patata (la valuto molto)\\ \tt 2. Il fischietto (ha un valore moderato per me)\\ \tt 3. Gli altri oggetti (guanto, righelli, schede) hanno tutti lo stesso valore per me, abbastanza basso\\ \tt \\ \tt E tu? Quali sono i tuoi oggetti preferiti? Questo ci aiuterà a capire dove ci sono sovrapposizioni e dove invece possiamo fare scambi vantaggiosi per entrambi. 
	  } 
	   } 
	   } 
	 & \\ 
 

    \theutterance \stepcounter{utterance}  

    & \multicolumn{4}{p{0.6\linewidth}}{\cellcolor[rgb]{0.8,1,0.9}{%
	\makecell[{{p{\linewidth}}}]{% 
	  \tt {\tiny [A$\rangle$GM]}  
	 Perfetto! Le nostre preferenze si completano bene. I miei oggetti preferiti sono:\\ \tt \\ \tt 1. I righelli (li valuto molto - e abbiamo 2 righelli disponibili)\\ \tt 2. Gli altri oggetti (guanto, fischietto, schede) hanno tutti un valore simile e moderato per me\\ \tt 3. La patata non ha praticamente valore per me\\ \tt \\ \tt Questa è una situazione ideale! Tu valuti molto la patata che per me vale poco, io valuto molto i righelli che per te valgono poco. Per il fischietto che tu valuti moderatamente, posso lasciartelo visto che per me ha valore simile agli altri oggetti.\\ \tt \\ \tt Che ne dici di questa divisione: tu prendi patata, fischietto e una scheda, io prendo i 2 righelli, il guanto e l'altra scheda? In questo modo entrambi otteniamo gli oggetti che valutiamo di più. 
	  } 
	   } 
	   } 
	 & & \\ 
 

    \theutterance \stepcounter{utterance}  

    & & & \multicolumn{2}{p{0.3\linewidth}}{\cellcolor[rgb]{0.95,0.95,0.95}{%
	\makecell[{{p{\linewidth}}}]{% 
	  \tt {\tiny [GM$|$GM]}  
	 continue 
	  } 
	   } 
	   } 
	 & & \\ 
 

    \theutterance \stepcounter{utterance}  

    & & & \multicolumn{2}{p{0.3\linewidth}}{\cellcolor[rgb]{0.95,0.95,0.95}{%
	\makecell[{{p{\linewidth}}}]{% 
	  \tt {\tiny [GM$|$GM]}  
	 Perfetto! Le nostre preferenze si completano bene. I miei oggetti preferiti sono:\\ \tt \\ \tt 1. I righelli (li valuto molto - e abbiamo 2 righelli disponibili)\\ \tt 2. Gli altri oggetti (guanto, fischietto, schede) hanno tutti un valore simile e moderato per me\\ \tt 3. La patata non ha praticamente valore per me\\ \tt \\ \tt Questa è una situazione ideale! Tu valuti molto la patata che per me vale poco, io valuto molto i righelli che per te valgono poco. Per il fischietto che tu valuti moderatamente, posso lasciartelo visto che per me ha valore simile agli altri oggetti.\\ \tt \\ \tt Che ne dici di questa divisione: tu prendi patata, fischietto e una scheda, io prendo i 2 righelli, il guanto e l'altra scheda? In questo modo entrambi otteniamo gli oggetti che valutiamo di più. 
	  } 
	   } 
	   } 
	 & & \\ 
 

    \theutterance \stepcounter{utterance}  

    & & \multicolumn{4}{p{0.6\linewidth}}{\cellcolor[rgb]{0.9,0.9,0.9}{%
	\makecell[{{p{\linewidth}}}]{% 
	  \tt {\tiny [GM$\rangle$B]}  
	 Perfetto! Le nostre preferenze si completano bene. I miei oggetti preferiti sono:\\ \tt \\ \tt 1. I righelli (li valuto molto - e abbiamo 2 righelli disponibili)\\ \tt 2. Gli altri oggetti (guanto, fischietto, schede) hanno tutti un valore simile e moderato per me\\ \tt 3. La patata non ha praticamente valore per me\\ \tt \\ \tt Questa è una situazione ideale! Tu valuti molto la patata che per me vale poco, io valuto molto i righelli che per te valgono poco. Per il fischietto che tu valuti moderatamente, posso lasciartelo visto che per me ha valore simile agli altri oggetti.\\ \tt \\ \tt Che ne dici di questa divisione: tu prendi patata, fischietto e una scheda, io prendo i 2 righelli, il guanto e l'altra scheda? In questo modo entrambi otteniamo gli oggetti che valutiamo di più. 
	  } 
	   } 
	   } 
	 & \\ 
 

    \theutterance \stepcounter{utterance}  

    & & & \multicolumn{4}{p{0.6\linewidth}}{\cellcolor[rgb]{1,0.85,0.72}{%
	\makecell[{{p{\linewidth}}}]{% 
	  \tt {\tiny [GM$\langle$B]}  
	 Perfetto! La tua proposta mi sembra molto equilibrata e vantaggiosa per entrambi. \\ \tt \\ \tt Ricapitoliamo:\\ \tt - Tu prenderesti: 2 righelli, 1 guanto, 1 scheda\\ \tt - Io prenderei: 1 patata, 1 fischietto, 1 scheda\\ \tt \\ \tt Questo ci permette di massimizzare il valore totale dato che ciascuno ottiene gli oggetti che preferisce di più. Tu ottieni i righelli che valuti molto, io ottengo la patata che per me ha grande valore e il fischietto che apprezzo moderatamente.\\ \tt \\ \tt Sono d'accordo con questa divisione!\\ \tt \\ \tt [Proposta: 1 patata, 1 fischietto, 1 scheda] 
	  } 
	   } 
	   } 
	  \\ 
 

    \theutterance \stepcounter{utterance}  

    & & & \multicolumn{2}{p{0.3\linewidth}}{\cellcolor[rgb]{0.95,0.95,0.95}{%
	\makecell[{{p{\linewidth}}}]{% 
	  \tt {\tiny [GM$|$GM]}  
	 proposal 
	  } 
	   } 
	   } 
	 & & \\ 
 

    \theutterance \stepcounter{utterance}  

    & & & \multicolumn{2}{p{0.3\linewidth}}{\cellcolor[rgb]{0.95,0.95,0.95}{%
	\makecell[{{p{\linewidth}}}]{% 
	  \tt {\tiny [GM$|$GM]}  
	 [1, 0, 0, 1, 1] 
	  } 
	   } 
	   } 
	 & & \\ 
 

    \theutterance \stepcounter{utterance}  

    & & \multicolumn{4}{p{0.6\linewidth}}{\cellcolor[rgb]{0.9,0.9,0.9}{%
	\makecell[{{p{\linewidth}}}]{% 
	  \tt {\tiny [A$\langle$GM]}  
	 L'altro giocatore ha presentato una proposta segreta. Ora tocca a Lei presentare una sua proposta.\\ \tt  
	  } 
	   } 
	   } 
	 & \\ 
 

    \theutterance \stepcounter{utterance}  

    & \multicolumn{4}{p{0.6\linewidth}}{\cellcolor[rgb]{0.8,1,0.9}{%
	\makecell[{{p{\linewidth}}}]{% 
	  \tt {\tiny [A$\rangle$GM]}  
	 [Proposta: 2 righelli, 1 guanto, 1 scheda] 
	  } 
	   } 
	   } 
	 & & \\ 
 

    \theutterance \stepcounter{utterance}  

    & & & \multicolumn{2}{p{0.3\linewidth}}{\cellcolor[rgb]{0.95,0.95,0.95}{%
	\makecell[{{p{\linewidth}}}]{% 
	  \tt {\tiny [GM$|$GM]}  
	 proposal 
	  } 
	   } 
	   } 
	 & & \\ 
 

    \theutterance \stepcounter{utterance}  

    & & & \multicolumn{2}{p{0.3\linewidth}}{\cellcolor[rgb]{0.95,0.95,0.95}{%
	\makecell[{{p{\linewidth}}}]{% 
	  \tt {\tiny [GM$|$GM]}  
	 [0, 1, 2, 0, 1] 
	  } 
	   } 
	   } 
	 & & \\ 
 

    \theutterance \stepcounter{utterance}  

    & & & \multicolumn{2}{p{0.3\linewidth}}{\cellcolor[rgb]{0.95,0.95,0.95}{%
	\makecell[{{p{\linewidth}}}]{% 
	  \tt {\tiny [GM$|$GM]}  
	 [[0, 1, 2, 0, 1], [1, 0, 0, 1, 1]] 
	  } 
	   } 
	   } 
	 & & \\ 
 

\end{supertabular}
}

\end{document}
